\subsection{对应的逆}

设 \(G\) 为一个图，其投影为 \(A = \operatorname{pr}_1 G\) , \(B = \operatorname{pr}_2 G\) 。关系 \((y, x) \in G\) 蕴含 \((x, y) \in B \times A\) ；因此该关系有一个图，它由所有有序对 \((x, y)\) （满足 \((y, x) \in G\) ）组成。

定义5. 设G为一个图。其元素为有序对 \((x, y)\) 使得 \((y, x) \in G\) 的图称为G的逆，并记为 \(\overline{G}\) .

对于每个集合X， \(\overline{G}\langle X\rangle\) 称为X在G下的逆像。

显然， \(\overline{G}\) 的逆是 \(G\) ，以及

\[
\mathrm{pr} _ {1} ^ {- 1} \mathrm{G} = \mathrm{pr} _ {2} \mathrm{G}, \quad \mathrm{pr} _ {2} ^ {- 1} \mathrm{G} = \mathrm{pr} _ {1} \mathrm{G}.
\]

特别地，如果 X、Y 是两个集合，我们有

\[
\widehat {\mathbf {X} \times \mathbf {Y}} = \mathbf {Y} \times \mathbf {X}.
\]

图 G 被称为对称的，如果 \(\overline{G}^{-1}=G\) .

让 \(\Gamma = (G, A, B)\) 是之间的对应关系 \(A\) 并且 \(B\) 。由于 \(\operatorname{pr}_1^{\overline{-1}} G \subset B\) 并且 \(\operatorname{pr}_2^{\overline{-1}} G \subset A\) ，三元组 \((G, B, A)\) 是...之间的对应关系 \(B\) 并且 \(A\) ，称为该对应的逆对应 \(\Gamma\) ，并记作 \(\overline{\Gamma}\) . 对于每一个子集 \(Y\) 的 \(B\) ，图像 \(\overline{\Gamma} \langle Y \rangle\) 的 \(Y\) 在……之下 \(\overline{\Gamma}\) 也被称为 \(Y\) 在……之下 \(\Gamma\) 。显然， \(\overline{\Gamma}\) 是 \(\Gamma\) .

